\scanpage{20}
\mspage{8}

We still have \deleted{to} see if $W\subset\widetilde W$. For the opposite inclusion we had to make restrictions on the $I$'s \inserted{(in \unsure{terms|favour} of $W$)} (condition 2), now to get $\widetilde W$ large enough, i.e. $TC$ large enough, we must make an assumption which will ensure that the bunch of $I$'s is large enough --- a natural condition would be that it contain a separating interval. We know already there \deleted{\unsure{are}}\inserted{\unsure{is one|are}} ``eligible'' \uline{separating} intervals, namely $W$-aspheric over $e$ and separating --- such as $L$.

% slanted
\margin{but this \unsure{cannot|can no} possibly be enough, \deleted{\unsure{it}} as it would (being fixed by itself) \deleted{\unsure{it}} give just one possible $\unsure{F'|F}$, hence one particular $\widetilde W$ being forced upon $W$....}

Let $f\in W$, consider factorization
\diagram{\begin{tikzcd}[ampersand replacement=\&,column sep=large]
X \arrow[r,"{\deleted{TC}\;i}"] \& X' \arrow[r,dash,"p"] \& Y \& i\in TC,\quad p\in F\quad\text{(or }i\in C,\;p\in TF\text{)}
\end{tikzcd}}
we know already that $i\in W$ (resp.! $p\in W$) \unsure{therefrom|therefore} $p\in W$ (resp. $i\in W$) --- hence we are reduced to \deleted{\unsure{the case $f$ is|Then one of}} proving either one of the two relations
\[
\begin{aligned}
 F\cap W&\subset TF &&\text{(necessary for }W\subset\widetilde W,\text{ as}\\
 &&& TF=F\cap\widetilde W\text{)}\\
 C\cap W&\subset TC &&(\mathrm{id})
\end{aligned}
\]
which are \inserted{\deleted{\unsure{both|not}}} equivalent \deleted{\unsure{each is|hence}} each to the condition $W\subset\widetilde W$ (or equivalently, $\deleted{\unsure{W=}}\widetilde W=W$). The second \inserted{[first?]} means that if $f:X-Y$ \inserted{[is in $W$ and]} has the \deleted{right} RLP with respect to $(TC)_0$ i.e. \inserted{to \deleted{\unsure{all}}} all $\alpha(\deleted{\unsure{i}}\unsure{i_1},i_0,I,e)$, then it has the RLP even with respect to \uline{all} monomorphisms which are $\in W$. \inserted{[????]}

We would like this to hold at least in the following two cases:
% slanted
\margin{In both cases $\mathscr A=A^{\wedge}$.}

a) $W=\underline W_A$ ($\underline W$ a basic localizer), $A$ a (local?) $\underline W$-test category

b) $A=A_0\times T$, \inserted{$\widehat A\simeq\underline{\mathrm{Hom}}(A_0^{\mathrm{opp}},T^\wedge)\simeq(A_0\unsure{\times}T)^\wedge$} with $T$ a test category, $W=\underline W_{A/A_0}$ ($\underline W$ a basic localizer, $\underline W_{A/A_0}$ is defined componentwise via $\underline W_T$.
% archive sheet number: 19
